% !TeX root = ../main.tex
\section*{Выводы}
В ходе выполнения работы была успешно смоделирована и исследована работа устройства цифровой квадратурной демодуляции радиосигналов. На основе результатов имитационного моделирования можно сделать следующие выводы:

\begin{enumerate}
    \item \textbf{Анализ спектральных преобразований:} Подтверждена математическая модель цифрового гетеродинирования. Наглядно продемонстрирован эффект алиасинга (наложения спектров), возникающий из-за намеренной субдискретизации входного радиосигнала, когда частота дискретизации АЦП меньше его несущей частоты.
    
    \item \textbf{Влияние разрядности цифрового гетеродина ($n_g$):} Установлено, что при использовании целочисленной арифметики (с фиксированной запятой) ошибки квантования формируют интенсивные паразитные нечетные гармоники (шпоры). Увеличение количества бит ЦГ снижает уровень шума квантования и подавляет эти гармоники. Для исключения влияния ЦГ на общий динамический диапазон приемника его разрядность должна превышать разрядность АЦП.
    
    \item \textbf{Влияние амплитудно-фазового рассогласования:} Исследование квадратурных каналов показало, что нарушение их строгой ортогональности приводит к появлению паразитной зеркальной составляющей в спектре выходного сигнала. Фазовая ошибка формирует выраженный зеркальный пик, а совместное наличие амплитудного и фазового дисбаланса максимизирует мощность этой мешающей компоненты, что критически ухудшает качество демодулированного сигнала.
\end{enumerate}